% GENERATED FILE. Edit canonical lessons, then run npm run generate:books.
% canonical-source-hash: fnv1a64:9d4db091729da4fe
% canonical-lessons: GE-C423-anker, GE-C423-ankern, GE-C423-leuchtturm, GE-C423-matrose, GE-C423-kentern

\chapter{Im Hafen}
\label{ch:ge-im-hafen}
\hlchaptermodality{423}

\begin{chapteropening}
\textbf{By the end of this chapter:} \emph{I can talk about a harbour: the anchor, lying at anchor, the lighthouse, the sailor and capsizing.}

Complete the last lesson of chapter 423: I can talk about a harbour: the anchor, lying at anchor, the lighthouse, the sailor and capsizing.
\end{chapteropening}

\section[der Anker]{der Anker — the anchor}

\label{lesson:GE-C423-anker}



\begin{quote}
Before the new one: say the German for to rock, then the German for the captain.
\end{quote}

\subsection*{You'll want to know: der Anker}
\textbf{der Anker} — \textquotedblleft{}the anchor\textquotedblright{}.

\begin{quote}
Der Anker fällt ins Wasser.
\end{quote}

\begin{quote}
\emph{The anchor drops into the water.}
\end{quote}

Same form in the plural. \emph{Vor Anker liegen} — to lie at anchor.

\begin{cousinweb}
\emph{Anker} is an early borrowing from Latin \emph{ancora} 'anchor' (Lewis \& Short), already in Old High German as \emph{anker} (EtymDB). English \emph{anchor} comes from the same Latin word.
\end{cousinweb}

\begin{grammarlens}[title={What you've built}]
One more word for telling, explaining and describing at B1.
\end{grammarlens}

\subsection*{Your turn}
\begin{itemize}
\raggedright
  \item \emph{Say it:} \emph{der Anker}
  \item \emph{Say it:} \emph{der Anker}, once more
  \item \emph{Say it:} a sentence of your own with \emph{der Anker}
  \item \emph{Recall:} say the German for to rock, then the German for the captain, then say \emph{der Anker} again
\end{itemize}

\begin{tcolorbox}[breakable,colback=teal!4,colframe=teal!35!black,title={Before you move on}]
What does der Anker mean? (\textquotedblleft{}The anchor\textquotedblright{}.) Say it once more.
\end{tcolorbox}
\section[ankern]{ankern — to anchor, to lie at anchor}

\label{lesson:GE-C423-ankern}



\begin{quote}
Before the new one: say the German for the captain, then the German for the anchor.
\end{quote}

\subsection*{You'll want to know: ankern}
\textbf{ankern} — \textquotedblleft{}to anchor, to lie at anchor\textquotedblright{}.

\begin{quote}
Das Boot ankert vor dem Strand.
\end{quote}

\begin{quote}
\emph{The boat is anchored off the beach.}
\end{quote}

\emph{Es ankert}, \emph{es hat geankert}.

\begin{cousinweb}
\emph{ankern} is built on \emph{Anker} 'anchor', an early borrowing from Latin \emph{ancora} (EtymDB; Lewis \& Short).
\end{cousinweb}

\begin{grammarlens}[title={What you've built}]
One more word for telling, explaining and describing at B1.
\end{grammarlens}

\subsection*{Your turn}
\begin{itemize}
\raggedright
  \item \emph{Say it:} \emph{ankern}
  \item \emph{Say it:} \emph{ankern}, once more
  \item \emph{Say it:} a sentence of your own with \emph{ankern}
  \item \emph{Recall:} say the German for the captain, then the German for the anchor, then say \emph{ankern} again
\end{itemize}

\begin{tcolorbox}[breakable,colback=teal!4,colframe=teal!35!black,title={Before you move on}]
What does ankern mean? (\textquotedblleft{}To anchor, to lie at anchor\textquotedblright{}.) Say it once more.
\end{tcolorbox}
\section[der Leuchtturm]{der Leuchtturm — the lighthouse}

\label{lesson:GE-C423-leuchtturm}



\begin{quote}
Before the new one: say the German for the anchor, then the German for to anchor.
\end{quote}

\subsection*{You'll want to know: der Leuchtturm}
\textbf{der Leuchtturm} — \textquotedblleft{}the lighthouse\textquotedblright{}.

\begin{quote}
Nachts sieht man das Licht vom Leuchtturm.
\end{quote}

\begin{quote}
\emph{At night the light from the lighthouse can be seen.}
\end{quote}

Plural: \emph{die Leuchttürme}.

\begin{cousinweb}
\emph{Leuchtturm} is built from \emph{leuchten} 'to shine', an inherited Germanic verb related to \emph{Licht} 'light', and \emph{Turm} 'tower', which goes back through Old High German \emph{turn} to Latin \emph{turris} (EtymDB).
\end{cousinweb}

\begin{grammarlens}[title={What you've built}]
One more word for telling, explaining and describing at B1.
\end{grammarlens}

\subsection*{Your turn}
\begin{itemize}
\raggedright
  \item \emph{Say it:} \emph{der Leuchtturm}
  \item \emph{Say it:} \emph{der Leuchtturm}, once more
  \item \emph{Say it:} a sentence of your own with \emph{der Leuchtturm}
  \item \emph{Recall:} say the German for the anchor, then the German for to anchor, then say \emph{der Leuchtturm} again
\end{itemize}

\begin{tcolorbox}[breakable,colback=teal!4,colframe=teal!35!black,title={Before you move on}]
What does der Leuchtturm mean? (\textquotedblleft{}The lighthouse\textquotedblright{}.) Say it once more.
\end{tcolorbox}
\section[der Matrose]{der Matrose — the sailor}

\label{lesson:GE-C423-matrose}



\begin{quote}
Before the new one: say the German for to anchor, then the German for the lighthouse.
\end{quote}

\subsection*{You'll want to know: der Matrose}
\textbf{der Matrose} — \textquotedblleft{}the sailor\textquotedblright{}.

\begin{quote}
Ein Matrose wirft das Seil an Land.
\end{quote}

\begin{quote}
\emph{A sailor throws the rope ashore.}
\end{quote}

Plural: \emph{die Matrosen}; weak noun: \emph{den Matrosen}.

\begin{cousinweb}
\emph{Matrose} was borrowed from Dutch \emph{matroos}, which came from French \emph{matelot} 'sailor', perhaps first 'bunkmate' (EtymDB).
\end{cousinweb}

\begin{grammarlens}[title={What you've built}]
One more word for telling, explaining and describing at B1.
\end{grammarlens}

\subsection*{Your turn}
\begin{itemize}
\raggedright
  \item \emph{Say it:} \emph{der Matrose}
  \item \emph{Say it:} \emph{der Matrose}, once more
  \item \emph{Say it:} a sentence of your own with \emph{der Matrose}
  \item \emph{Recall:} say the German for to anchor, then the German for the lighthouse, then say \emph{der Matrose} again
\end{itemize}

\begin{tcolorbox}[breakable,colback=teal!4,colframe=teal!35!black,title={Before you move on}]
What does der Matrose mean? (\textquotedblleft{}The sailor\textquotedblright{}.) Say it once more.
\end{tcolorbox}
\section[kentern]{kentern — to capsize}

\label{lesson:GE-C423-kentern}



\begin{quote}
Before the new one: say the German for the lighthouse, then the German for the sailor.
\end{quote}

\subsection*{You'll want to know: kentern}
\textbf{kentern} — \textquotedblleft{}to capsize\textquotedblright{}.

\begin{quote}
Bei starkem Wind ist ein kleines Boot gekentert.
\end{quote}

\begin{quote}
\emph{In a strong wind a small boat capsized.}
\end{quote}

Perfekt with \emph{sein}: \emph{das Boot ist gekentert}.

\begin{cousinweb}
\emph{kentern} came into standard German from Low German \emph{kenteren}, a sailors' word; it is usually linked to \emph{Kante} 'edge', a boat tipped over on its side. EtymDB does not list it, so the line is hedged.
\end{cousinweb}

\begin{grammarlens}[title={What you've built}]
One more word for telling, explaining and describing at B1.
\end{grammarlens}

\subsection*{Your turn}
\begin{itemize}
\raggedright
  \item \emph{Say it:} \emph{kentern}
  \item \emph{Say it:} \emph{kentern}, once more
  \item \emph{Say it:} a sentence of your own with \emph{kentern}
  \item \emph{Recall:} say the German for the lighthouse, then the German for the sailor, then say \emph{kentern} again
\end{itemize}

\begin{tcolorbox}[breakable,colback=teal!4,colframe=teal!35!black,title={Before you move on}]
What does kentern mean? (\textquotedblleft{}To capsize\textquotedblright{}.) Say it once more.
\end{tcolorbox}
